\chapter{From Data to Conjecture}
\label{ch:experiments}

\begin{center}\itshape
The computer does not prove theorems. It tells you which theorems to try to prove,\\
and, more often, which ones not to waste a year on.
\end{center}

In 1995 a computer at Simon Fraser University in Vancouver found a new formula for $\pi$. Simon Plouffe, David Bailey and Peter Borwein had been
looking for linear relations between $\pi$ and a handful of slowly converging series, with an integer relation algorithm running at high
precision, and one day it returned a vector of small integers:
\[
\pi=\sum_{k=0}^\infty\frac1{16^k}\Bigl(\frac4{8k+1}-\frac2{8k+4}-\frac1{8k+5}-\frac1{8k+6}\Bigr).
\]
The formula is not faster than the classical ones. Its point is elsewhere: because of the $16^k$, it allows one to compute the millionth
hexadecimal digit of $\pi$ without computing the first $999\,999$. Nobody had suspected that this was possible. And once the computer had
found the formula, its proof took a few lines of first-year calculus.

This chapter is about that sequence of events: compute, notice, guess, test, prove. We prove the formula and the digit-extraction algorithm,
then the theorem behind the integer relation search, the lattice reduction algorithm of Lenstra, Lenstra and Lov\'asz; and we end with the
other side of the coin: patterns that hold for a billion cases and then fail.

\section{The formula and its proof}

\begin{theorem}[Bailey, Borwein and Plouffe, 1995]\label{thm:bbp}
For $j\ge1$, let $S_j=\sum_{k\ge0}\frac1{16^k(8k+j)}$. Then $\pi=4S_1-2S_4-S_5-S_6$.
\end{theorem}

\begin{proof}
For $k\ge0$, $\int_0^{1/\sqrt2}x^{8k+j-1}\dd x=\frac{2^{-j/2}}{16^k(8k+j)}$. Summing over $k$ (the terms are positive, so the sum and the integral can be exchanged),
\[
S_j=2^{j/2}\int_0^{1/\sqrt2}\frac{x^{j-1}}{1-x^8}\dd x.
\]
So
\[
4S_1-2S_4-S_5-S_6=\int_0^{1/\sqrt2}\frac{4\sqrt2-8x^3-4\sqrt2x^4-8x^5}{1-x^8}\dd x,
\]
and the substitution $x=y/\sqrt2$ turns this into
\[
\int_0^1\frac{16(4-2y^3-y^4-y^5)}{16-y^8}\dd y=\int_0^1\frac{16(y-1)}{(y^2-2)(y^2-2y+2)}\dd y,
\]
since $16-y^8=(4-y^4)(4+y^4)$, $4-y^4=-(y^2-2)(y^2+2)$, $4+y^4=(y^2-2y+2)(y^2+2y+2)$, and $4-2y^3-y^4-y^5=(1-y)(y^2+2)(y^2+2y+2)$, as one checks by expanding. In partial fractions,
\[
\frac{16(y-1)}{(y^2-2)(y^2-2y+2)}=\frac{4y}{y^2-2}-\frac{4(y-1)}{(y-1)^2+1}+\frac4{(y-1)^2+1},
\]
(multiply out: the numerator is $4y(y^2-2y+2)-(4y-8)(y^2-2)=16y-16$). The three integrals over $[0,1]$ are $2\ln\frac12$, $-2\ln\frac12$ and $4\bigl(\arctan0-\arctan(-1)\bigr)=\pi$. Their sum is $\pi$.
\end{proof}

The proof is easy \emph{once one knows what to prove}. That is the typical shape of an experimental discovery: the hard part, finding the
statement, was done by a machine, and the human part, the proof, was routine. The opposite also happens; we shall meet it in Chapter~\ref{ch:heuristics}.

\begin{proposition}[digit extraction]\label{prop:digits}
For $d\ge0$ and $j\ge1$, let
\[
T_j(d)=\sum_{k=0}^d\frac{16^{d-k}\bmod(8k+j)}{8k+j}+\sum_{k>d}\frac{16^{d-k}}{8k+j}.
\]
Then the fractional part of $16^d\pi$ equals the fractional part of $4T_1(d)-2T_4(d)-T_5(d)-T_6(d)$. In particular the hexadecimal digits of $\pi$ from position $d+1$ on can be
computed with $O(d\log d)$ arithmetic operations on numbers of size $O(\log d)$ bits, without computing the previous digits.
\end{proposition}

\begin{proof}
For $k\le d$, $\frac{16^{d-k}}{8k+j}$ differs from $\frac{16^{d-k}\bmod(8k+j)}{8k+j}$ by an integer. So $16^dS_j-T_j(d)$ is an integer, and as the coefficients $4,-2,-1,-1$ are integers,
$16^d\pi-(4T_1-2T_4-T_5-T_6)$ is an integer by Theorem~\ref{thm:bbp}. Each residue $16^{d-k}\bmod(8k+j)$ is computed by repeated squaring modulo $8k+j$ with $O(\log d)$ multiplications of
numbers below $8d+8$, and there are $O(d)$ of them; the tail is less than $\sum_{m\ge1}16^{-m}=\frac1{15}$ and a few terms give it to any fixed accuracy.
\end{proof}

In floating point, each $T_j(d)$ is a sum of $d$ numbers in $[0,1)$, computed with relative error about $10^{-16}$; the error in the final fractional part
is then about $d\cdot10^{-16}$, so for $d=10^6$ one gets about eight correct hexadecimal digits. Direct computation gives the digits of $\pi$ starting at position $10^6$ in a few seconds:
\texttt{26C65E52}\dots, the same digits that Bailey, Borwein and Plouffe printed in 1997.

\begin{think}
Is there a formula of this kind for the \emph{decimal} digits of $\pi$, that is, with $10^k$ in place of $16^k$? Nobody knows one. What would the
digit-extraction argument need from such a formula?
\end{think}

\section{How the formula was found: lattices}

An \emph{integer relation} between real numbers $x_1,\dots,x_n$ is a nonzero $m\in\Z^n$ with $m_1x_1+\dots+m_nx_n=0$. Theorem~\ref{thm:bbp} is the relation
$(1,-4,2,1,1)$ between $(\pi,S_1,S_4,S_5,S_6)$. One may search among $\pi$ and $S_1,\dots,S_8$; with the numbers known to a few hundred digits, how does one
find a relation with small coefficients, if there is one? The answer is a short vector in a lattice.

Let $b_1,\dots,b_n$ be linearly independent vectors of $\R^m$. The \emph{lattice} they generate is $L=\Z b_1+\dots+\Z b_n$, and $\lambda_1(L)$ is the length of its shortest nonzero vector.
Finding a shortest vector is hard in general, but finding one that is short \emph{up to a factor $2^{(n-1)/2}$} is easy, and that is enough for many purposes.
Write $b_i^*$ for the Gram--Schmidt vectors: $b_1^*=b_1$, and $b_i^*$ is the component of $b_i$ orthogonal to $b_1,\dots,b_{i-1}$, so that
\[
b_i=b_i^*+\sum_{j<i}\mu_{ij}b_j^*,\qquad\mu_{ij}=\frac{\langle b_i,b_j^*\rangle}{|b_j^*|^2}.
\]

\begin{definition}\label{def:lll}
The basis is \emph{LLL-reduced} if it is \emph{size-reduced}, $|\mu_{ij}|\le\frac12$ for all $j<i$, and satisfies the \emph{Lov\'asz condition}
\[
|b_i^*|^2\ge\Bigl(\frac34-\mu_{i,i-1}^2\Bigr)|b_{i-1}^*|^2\qquad(2\le i\le n).
\]
\end{definition}

\begin{lemma}\label{lem:gsbound}
Every nonzero $v\in L$ satisfies $|v|\ge\min_i|b_i^*|$.
\end{lemma}

\begin{proof}
Write $v=\sum c_ib_i$ with $c_i\in\Z$, and let $k$ be the largest index with $c_k\ne0$. Expanding each $b_i$ in the Gram--Schmidt basis, $v=c_kb_k^*+\sum_{j<k}\nu_jb_j^*$ for some reals $\nu_j$. The $b_j^*$ are orthogonal, so
$|v|\ge|c_k|\,|b_k^*|\ge|b_k^*|$.
\end{proof}

\begin{theorem}[Lenstra, Lenstra and Lov\'asz, 1982]\label{thm:lll}
Let $b_1,\dots,b_n$ be a basis of the lattice $L$.
\begin{enumerate}[label=(\alph*)]
\item If $b_1,\dots,b_n$ is LLL-reduced, then $|b_1|\le2^{(n-1)/2}\lambda_1(L)$.
\item If $b_1,\dots,b_n\in\Z^m$, the algorithm below stops, and returns an LLL-reduced basis of the same lattice, after at most $n+2\log_{4/3}D_0$ passes through its loop, where, for the input basis,
\[
D_0=\prod_{i<n}\det\bigl(\langle b_k,b_l\rangle\bigr)_{k,l\le i}.
\]
\end{enumerate}
\end{theorem}

\emph{The algorithm.} Set $k=2$. While $k\le n$: size-reduce $b_k$, that is, for $j=k-1,k-2,\dots,1$ replace $b_k$ by $b_k-\lfloor\mu_{kj}\rceil b_j$ (nearest integer); then if the Lov\'asz condition
holds at $k$, increase $k$ by $1$; otherwise exchange $b_{k-1}$ and $b_k$, and set $k=\max(k-1,2)$.

\begin{proof}
(a) Size reduction and the Lov\'asz condition give $|b_i^*|^2\ge(\frac34-\frac14)|b_{i-1}^*|^2=\frac12|b_{i-1}^*|^2$, so $|b_1|^2=|b_1^*|^2\le2^{i-1}|b_i^*|^2\le2^{n-1}|b_i^*|^2$ for every $i$. With
Lemma~\ref{lem:gsbound}, $|b_1|^2\le2^{n-1}\min_i|b_i^*|^2\le2^{n-1}\lambda_1(L)^2$.

(b) \emph{The lattice does not change}: both operations are invertible integer changes of basis.

\emph{Invariant: at the start of each pass, $b_1,\dots,b_{k-1}$ is LLL-reduced.} It holds for $k=2$. Replacing $b_k$ by $b_k-qb_j$ with $j<k$ does not change $b_1,\dots,b_{k-1}$ nor any $b_i^*$, and changes
$\mu_{kl}$ only for $l\le j$, into $\mu_{kl}-q\mu_{jl}$ (with $\mu_{jj}=1$); so after the loop over $j=k-1,\dots,1$, each $\mu_{kj}$ was last changed when it was made at most $\frac12$ in absolute value. If the Lov\'asz
condition then holds at $k$, $b_1,\dots,b_k$ is LLL-reduced, and $k$ increases. Otherwise the exchange leaves $b_1,\dots,b_{k-2}$ untouched, and $k$ becomes $\max(k-1,2)$; the invariant holds again. So when the
loop ends, with $k=n+1$, the basis is LLL-reduced.

\emph{Termination.} Let $d_i=\det(\langle b_k,b_l\rangle)_{k,l\le i}$, a Gram determinant. Gram--Schmidt is a change of basis with a unipotent triangular matrix, so $d_i=\prod_{j\le i}|b_j^*|^2>0$, and $d_i$
is an integer, as the $b_k$ have integer coordinates; so $D=\prod_{i<n}d_i$ is a positive integer. Size reduction does not change the $b_j^*$, hence not $D$. An exchange at $k$ changes only
$b_{k-1}^*$ and $b_k^*$ among the Gram--Schmidt vectors, and among the $d_i$ only $d_{k-1}$ (since $d_k$ is the Gram determinant of the same set of vectors). The new $b_{k-1}^*$ is the component of the old $b_k$
orthogonal to $b_1,\dots,b_{k-2}$, which is $b_k^*+\mu_{k,k-1}b_{k-1}^*$, of squared length $|b_k^*|^2+\mu_{k,k-1}^2|b_{k-1}^*|^2<\frac34|b_{k-1}^*|^2$ because the Lov\'asz condition failed. So $d_{k-1}$, and $D$, are multiplied
by less than $\frac34$. A positive integer can be multiplied by less than $\frac34$ at most $\log_{4/3}D_0$ times, so there are at most that many exchanges. Each pass either increases $k$ or
is an exchange that decreases $k$ by at most $1$; since $2\le k\le n+1$, the number of passes is at most $n-1+2\times(\text{number of exchanges})$.
\end{proof}

\begin{figure}[ht]
\centering
\fig{ch07-lattice}
\caption{The lattice generated by $b_1=(13,9)$ and $b_2=(5,4)$ (orange) is also generated by $b_1'=(-1,2)$ and $b_2'=(3,1)$ (blue). The first basis is long and nearly flat, the second
short and nearly orthogonal. Lattice reduction finds the second from the first; in dimension $2$ it is Lagrange's and Gauss's reduction of binary quadratic forms (Exercise~\ref{ex:gauss2}).}
\label{fig:lattice}
\end{figure}

With more care, one shows that the algorithm runs in time polynomial in $n$, $m$ and the number of bits of the input, since the numbers it manipulates stay of
polynomial size; we do not need this. The exponential factor $2^{(n-1)/2}$ in (a) looks terrible, and in the worst case it is nearly attained; in practice
LLL usually finds vectors much shorter than its guarantee.

\subsection*{From lattices to integer relations}

Let $x_1,\dots,x_n$ be real numbers and $N$ a large integer. Consider the lattice in $\Z^{n+1}$ with basis
\[
b_i=\bigl(e_i,\ \lfloor Nx_i\rceil\bigr)\qquad(1\le i\le n),
\]
where $e_i$ is the $i$-th unit vector of $\Z^n$ and $\lfloor\cdot\rceil$ is the nearest integer. A lattice vector is $\sum m_ib_i=(m,t)$ with $t=\sum m_i\lfloor Nx_i\rceil$, and $t\approx N\sum m_ix_i$. So a short vector has a small
$m$ and a small $\sum m_ix_i$: an approximate integer relation of small height.

\begin{proposition}[what the search guarantees]\label{prop:relation}
Suppose $x$ has an integer relation $m^\circ$. Let $(m,t)$ be the first vector of an LLL-reduced basis of this lattice. Then
\[
|m|\le C_n|m^\circ|\qquad\text{and}\qquad\Bigl|\sum_im_ix_i\Bigr|\le\frac{C_n'|m^\circ|}N,
\]
with $C_n=2^{(n-1)/2}\sqrt{1+n/4}$ and $C'_n=C_n(1+\frac{\sqrt n}2)$.
\end{proposition}

\begin{proof}
For any $m'\in\Z^n$, $|\sum m'_i\lfloor Nx_i\rceil-N\sum m'_ix_i|\le\frac12\sum|m'_i|\le\frac{\sqrt n}2|m'|$. For $m'=m^\circ$, this shows that $(m^\circ,t^\circ)$ is a lattice vector of length at most
$|m^\circ|\sqrt{1+n/4}$. By Theorem~\ref{thm:lll}(a), $|(m,t)|\le C_n|m^\circ|$, so $|m|\le C_n|m^\circ|$ and $|t|\le C_n|m^\circ|$; and $N|\sum m_ix_i|\le|t|+\frac{\sqrt n}2|m|$.
\end{proof}

The proposition does \emph{not} say that $m$ is a relation: only that it is a relation up to an error of order $1/N$, with a height comparable to that of the true one.
If the computation is done with $D$ digits and $N=10^D$, and if $\sum m_ix_i$ is zero to all $D$ digits while $m$ has three-digit entries, the evidence is overwhelming, but it is not a
proof. For BBP, the proof came afterwards, by calculus. In practice one uses PSLQ, an algorithm of Helaman Ferguson and David Bailey designed for exactly
this problem, which also certifies that no relation of height below a computed bound exists; LLL shows the principle.

\section{Coincidences}

The computer finds patterns; it also finds patterns that are not there. Three famous examples, in increasing order of danger.

\emph{Almost integers.} In April 1975, in his \emph{Scientific American} column, Martin Gardner announced that Ramanujan had conjectured, and a computer had now
confirmed, that $e^{\pi\sqrt{163}}$ is an integer. It was an April fool:
\[
e^{\pi\sqrt{163}}=262\,537\,412\,640\,768\,743.999\,999\,999\,999\,25\dots
\]
The near miss has a reason. By Tome~12, Theorem~14.3 and Remark~14.5, the modular invariant $j$ takes the integer value $-640320^3$ at $\tau=\frac{1+\sqrt{-163}}2$, and $j(\tau)=q^{-1}+744+196884q+\dots$
with $q=e^{2\pi i\tau}=-e^{-\pi\sqrt{163}}$. So $e^{\pi\sqrt{163}}=640320^3+744-196884e^{-\pi\sqrt{163}}+\dots$, and the error is about $7.5\times10^{-13}$ (Exercise~\ref{ex:ramanujan}). Here the coincidence is a theorem in
disguise; the lesson is that an almost integer deserves an explanation.

\emph{The Pólya conjecture.} Let $\lambda(n)=(-1)^{\Omega(n)}$, where $\Omega(n)$ counts the prime factors of $n$ with multiplicity, and $L(x)=\sum_{n\le x}\lambda(n)$. In 1919 George Pólya
observed that $L(x)\le0$ for all $x\ge2$ in his tables, which suggests that at least half the numbers up to $x$ have an odd number of prime factors. In 1958 C.~Brian Haselgrove proved that
$L(x)>0$ for some $x$, without finding one; in 1980 Minoru Tanaka found the smallest, $x=906\,150\,257$. Nine hundred million cases of evidence, and the conjecture was false.

\emph{Skewes.} Every table of primes ever printed shows $\pi(x)<\operatorname{li}(x)$, where $\operatorname{li}(x)$ is the logarithmic integral, essentially $\int_2^x\frac{\dd t}{\ln t}$. In 1914
John Edensor Littlewood proved that $\pi(x)-\operatorname{li}(x)$ changes sign infinitely often. The first sign change is beyond every computer search, and the best
current bounds place it near $10^{316}$.

Why do these fail? In both cases the quantity is governed by the zeros of the zeta function, and its sign is decided by a sum of oscillating terms of size
$x^{1/2}/\ln x$ against a bias of comparable size; the oscillations win only when $\ln x$ is large enough for the terms to align, and $\ln(10^9)\approx21$ is not large.
Richard Guy collected dozens of such examples in an essay titled \emph{The strong law of small numbers}: there are not enough small numbers to meet the many demands
made of them. The working rule: \emph{numerical evidence is evidence about the numbers you computed. Before you trust it, ask on what scale the phenomenon lives.} If it
lives on the scale $\ln\ln x$, a billion is small.

\begin{beauty}
A formula for $\pi$ that hides inside it an algorithm for its digits; an integer, $640320^3+744$, that a modular function forces to be within $10^{-12}$ of $e^{\pi\sqrt{163}}$; and a
lattice algorithm whose only job is to find short vectors, which became the key to factoring polynomials, breaking knapsack cryptosystems and, now, building the cryptography
meant to survive quantum computers. Experimental mathematics is not a substitute for proof; it is the art of asking the right question.
\end{beauty}

\section{Exercises}

\begin{exercise}[$\star$]\label{ex:log2}
Show that $\ln2=\sum_{k\ge1}\frac1{k2^k}$, and deduce an algorithm computing the binary digits of $\ln2$ from position $d+1$ on, without computing the previous ones.
\end{exercise}

\begin{exercise}[$\star$]\label{ex:bbpcheck}
Check numerically that the first terms of Theorem~\ref{thm:bbp} give $\pi$ to about $1.2$ decimal digits per term. Why $1.2$?
\end{exercise}

\begin{exercise}[$\star\star$]\label{ex:gauss2}
For $n=2$, show that if $|\mu_{21}|\le\frac12$ and $|b_2|\ge|b_1|$, then $b_1$ is a shortest nonzero vector of $L$. (This is the reduction of Lagrange and Gauss; the Lov\'asz condition with $\frac34$
replaced by $1$ says exactly $|b_2|\ge|b_1|$ when $\mu_{21}$ is taken into account.)
\end{exercise}

\begin{exercise}[$\star\star$]\label{ex:delta}
Replace $\frac34$ by $\delta\in(\frac14,1)$ in Definition~\ref{def:lll} and in the algorithm. Show that the algorithm still terminates and that $|b_1|\le(\delta-\frac14)^{-(n-1)/2}\lambda_1(L)$. What goes wrong in the proof of termination
for $\delta=1$?
\end{exercise}

\begin{exercise}[$\star\star$]\label{ex:minpoly}
Let $\alpha=\sqrt2+\sqrt3$. Use Proposition~\ref{prop:relation} with $x=(1,\alpha,\alpha^2,\alpha^3,\alpha^4)$ to guess an integer polynomial of degree $4$ vanishing at $\alpha$ (or compute it by hand), and then prove
that it vanishes at $\alpha$ and is irreducible over $\Q$.
\end{exercise}

\begin{exercise}[$\star\star$]\label{ex:ramanujan}
Using $j(\tau)=q^{-1}+744+196884q+O(|q|^2)$ and $j\bigl(\frac{1+\sqrt{-163}}2\bigr)=-640320^3$, show that $640320^3+744-e^{\pi\sqrt{163}}=196884e^{-\pi\sqrt{163}}+O(e^{-2\pi\sqrt{163}})$ and evaluate it.
\end{exercise}

\begin{exercise}[$\star\star\star$]\label{ex:polya}
\emph{Pólya's conjecture, the easy direction.} Show that $\sum_{n\le x}\lambda(n)\lfloor x/n\rfloor$ equals the number of squares up to $x$, that is $\lfloor\sqrt x\rfloor$. (Use $\sum_{d\mid n}\lambda(d)=1$ if $n$ is a square, $0$ otherwise.)
Explain why this identity makes a negative bias of $L(x)$ plausible, without proving it.
\end{exercise}

\section{Hints}

\begin{hint}{ex:log2}
$-\ln(1-x)=\sum x^k/k$ at $x=\frac12$. Then imitate Proposition~\ref{prop:digits} with $2^{d-k}\bmod k$.
\end{hint}

\begin{hint}{ex:bbpcheck}
The terms decrease like $16^{-k}$, and $\log_{10}16\approx1.2$.
\end{hint}

\begin{hint}{ex:gauss2}
Write a lattice vector as $c_1b_1+c_2b_2$ and use $|c_1b_1+c_2b_2|^2=(c_1+c_2\mu_{21})^2|b_1|^2+c_2^2|b_2^*|^2$.
\end{hint}

\begin{hint}{ex:delta}
Redo both proofs with $\delta$; for termination, the factor is now $\delta$.
\end{hint}

\begin{hint}{ex:minpoly}
$\alpha^2=5+2\sqrt6$.
\end{hint}

\begin{hint}{ex:ramanujan}
$q=e^{\pi i(1+\sqrt{-163})}=-e^{-\pi\sqrt{163}}$.
\end{hint}

\begin{hint}{ex:polya}
$\sum_{n\le x}\lambda(n)\lfloor x/n\rfloor=\sum_{m\le x}\sum_{d\mid m}\lambda(d)$.
\end{hint}

\section{Solutions}

\begin{solution}{ex:log2}
\idea The same trick as for $\pi$, with base $2$.
\writeup For $|x|<1$, $-\ln(1-x)=\sum_{k\ge1}\frac{x^k}k$; at $x=\frac12$, $\ln2=\sum\frac1{k2^k}$. Then $2^d\ln2=\sum_{k\le d}\frac{2^{d-k}}k+\sum_{k>d}\frac{2^{d-k}}k$, and modulo $1$ the first sum equals $\sum_{k\le d}\frac{2^{d-k}\bmod k}k$, whose terms
are computed by repeated squaring modulo $k$. The tail is less than $\frac1{d+1}$ and is summed directly. The fractional part of the result is the fractional part of $2^d\ln2$, whose binary expansion
is the digits of $\ln2$ from position $d+1$ on. The BBP paper discusses this case as the simplest example of the method.
\end{solution}

\begin{solution}{ex:bbpcheck}
\idea The tail after $K$ terms is a geometric series.
\writeup The $k$-th term is $16^{-k}$ times a quantity of order $\frac1{2k}$, so the tail after $K$ terms is $O(16^{-K}/K)$: each term adds $\log_{10}16\approx1.204$ correct decimal digits. With $K=10$ terms one gets
about $12$ digits; numerical evaluation confirms it.
\end{solution}

\begin{solution}{ex:gauss2}
\idea Expand the squared length in the orthogonal basis $b_1^*=b_1$, $b_2^*$.
\writeup Let $v=c_1b_1+c_2b_2\ne0$, with $c_i\in\Z$. As $b_2=b_2^*+\mu_{21}b_1$, $|v|^2=(c_1+c_2\mu_{21})^2|b_1|^2+c_2^2|b_2^*|^2$. If $c_2=0$, then $|v|=|c_1|\,|b_1|\ge|b_1|$. If $|c_2|\ge2$, then $|v|^2\ge4|b_2^*|^2$, and
$|b_2^*|^2=|b_2|^2-\mu_{21}^2|b_1|^2\ge(1-\frac14)|b_1|^2$, so $|v|^2\ge3|b_1|^2$. If $c_2=\pm1$, then $|v|^2=(c_1\pm\mu_{21})^2|b_1|^2+|b_2^*|^2$, and $|c_1\pm\mu_{21}|\ge|\mu_{21}|$ for every integer $c_1$ (with equality for $c_1=0$; otherwise $|c_1\pm\mu_{21}|\ge1-\frac12\ge|\mu_{21}|$); so
$|v|^2\ge\mu_{21}^2|b_1|^2+|b_2^*|^2=|b_2|^2\ge|b_1|^2$. So $|v|\ge|b_1|$ in all cases.
\end{solution}

\begin{solution}{ex:delta}
\idea Every $\frac34$ in the proofs becomes $\delta$.
\writeup With $\delta$, the Lov\'asz condition and size reduction give $|b_i^*|^2\ge(\delta-\frac14)|b_{i-1}^*|^2$, so $|b_1|^2\le(\delta-\frac14)^{-(i-1)}|b_i^*|^2$, and Lemma~\ref{lem:gsbound} gives $|b_1|\le(\delta-\frac14)^{-(n-1)/2}\lambda_1(L)$. For termination, an exchange
multiplies $D$ by less than $\delta$, and a positive integer can be so multiplied at most $\log_{1/\delta}D_0$ times. For $\delta=1$, an exchange only guarantees that $D$ \emph{decreases}; an integer
can decrease only finitely many times, so the algorithm still terminates, but the bound $\log_{1/\delta}D_0$ becomes infinite: no polynomial bound on the number of steps follows from this argument.
\end{solution}

\begin{solution}{ex:minpoly}
\idea Guess with the lattice, prove by algebra.
\writeup The search returns (up to sign) the vector $(1,0,-10,0,1)$, that is $P(x)=x^4-10x^2+1$. Proof that $P(\alpha)=0$: $\alpha^2=5+2\sqrt6$, so $(\alpha^2-5)^2=24$, $\alpha^4-10\alpha^2+1=0$. Irreducibility: the roots of $P$ are
$\pm\sqrt2\pm\sqrt3$, none rational, so $P$ has no linear factor over $\Q$. A monic quadratic factor over $\Q$ would have two of the roots as its roots, with rational sum and product. Let one of them be
$\sqrt2+\sqrt3$ (otherwise use the other factor). Its partner is $-\sqrt2-\sqrt3$ (product $-(5+2\sqrt6)\notin\Q$), $\sqrt2-\sqrt3$ (sum $2\sqrt2\notin\Q$) or $-\sqrt2+\sqrt3$ (sum $2\sqrt3\notin\Q$). So $P$ is irreducible over $\Q$.
\end{solution}

\begin{solution}{ex:ramanujan}
\idea Plug $q$ into the expansion of $j$.
\writeup At $\tau=\frac{1+\sqrt{-163}}2$, $q=e^{2\pi i\tau}=e^{\pi i}e^{-\pi\sqrt{163}}=-e^{-\pi\sqrt{163}}$, so $q^{-1}=-e^{\pi\sqrt{163}}$. The expansion gives $-640320^3=-e^{\pi\sqrt{163}}+744-196884e^{-\pi\sqrt{163}}+O(e^{-2\pi\sqrt{163}})$,
that is $640320^3+744-e^{\pi\sqrt{163}}=196884e^{-\pi\sqrt{163}}+O(e^{-2\pi\sqrt{163}})$. With $e^{\pi\sqrt{163}}\approx2.625\times10^{17}$, the difference is about $7.50\times10^{-13}$, as the decimal expansion shows.
\end{solution}

\begin{solution}{ex:polya}
\idea Count the pairs $(d,m)$ with $d\mid m\le x$ in two ways.
\writeup $\lambda$ is completely multiplicative, so $g=\sum_{d\mid n}\lambda(d)$ is multiplicative, with $g(p^a)=1-1+\dots+(-1)^a$, which is $1$ for even $a$ and $0$ for odd $a$; so $g(n)=1$ if $n$ is a square and $0$ otherwise.
Then $\sum_{n\le x}\lambda(n)\lfloor\frac xn\rfloor=\sum_{n\le x}\lambda(n)\#\{m\le x:n\mid m\}=\sum_{m\le x}g(m)=\lfloor\sqrt x\rfloor$. Heuristically, $\sum_n\lambda(n)\frac xn$ is then of order $\sqrt x$, which is much smaller than $x$: the
values of $\lambda$ must nearly cancel, but the weights $\frac1n$ favour small $n$, and the small $n$ (the primes, with $\lambda=-1$) are many. That pushes $L(x)$ down on average; Haselgrove's theorem says it
does not keep it down forever.
\end{solution}

\section*{Sources and further reading}
\addcontentsline{toc}{section}{Sources and further reading}

\begin{itemize}
\item \emph{BBP.} D.~Bailey, P.~Borwein and S.~Plouffe, ``On the rapid computation of various polylogarithmic constants'', \emph{Math.\ Comp.}\ 66 (1997), 903--913 (the formula, its discovery
in 1995, the digits at position $10^6$, and the $\ln2$ example). J.~Borwein and D.~Bailey, \emph{Mathematics by Experiment}, A~K Peters, 2004, for the experimental method in general.
\item \emph{Lattices.} A.~K.~Lenstra, H.~W.~Lenstra Jr.\ and L.~Lov\'asz, ``Factoring polynomials with rational coefficients'', \emph{Math.\ Ann.}\ 261 (1982), 515--534. H.~R.~P.~Ferguson, D.~H.~Bailey
and S.~Arno, ``Analysis of PSLQ, an integer relation finding algorithm'', \emph{Math.\ Comp.}\ 68 (1999), 351--369. P.~Q.~Nguyen and B.~Vall\'ee (eds.), \emph{The LLL Algorithm}, Springer, 2010.
\item \emph{Coincidences.} M.~Gardner, ``Six sensational discoveries that somehow or another have escaped public attention'', \emph{Scientific American} 232, no.~4 (April 1975), 127--131.
G.~P\'olya, ``Verschiedene Bemerkungen zur Zahlentheorie'', \emph{Jahresber.\ Deutsch.\ Math.-Verein.}\ 28 (1919), 31--40. C.~B.~Haselgrove, ``A disproof of a conjecture of P\'olya'', \emph{Mathematika} 5 (1958),
141--145. M.~Tanaka, ``A numerical investigation on cumulative sum of the Liouville function'', \emph{Tokyo J.\ Math.}\ 3 (1980), 187--189. J.~E.~Littlewood, ``Sur la distribution des nombres premiers'',
\emph{C.\ R.\ Acad.\ Sci.\ Paris} 158 (1914), 1869--1872. R.~K.~Guy, ``The strong law of small numbers'', \emph{Amer.\ Math.\ Monthly} 95 (1988), 697--712.
\end{itemize}
